\clearpage\NoSourceHead\SourceLeaf{88}
\renewcommand{\thefootnote}{\arabic{footnote})}
\setcounter{footnote}{0}
\begin{flushright}
\begin{minipage}{.70\textwidth}\small
»Dem Einzelnen bleibe die Freiheit, sich mit dem zu
beschäftigen, was ihn anzieht, was ihm Freude macht, was
ihm nützlich deucht; aber das eigentliche Studium der
Menschheit ist der Mensch.«\par
\begin{flushright}
\textsc{Goethe}, \textit{Wahlverwandtschaften}, VII. Kap.\par
Aus \textsc{Ottiliens} Tagebuch.
\end{flushright}
\end{minipage}
\end{flushright}
\bigskip
\begin{center}\bfseries I. Einleitendes und Übersicht über das benutzte Material.\end{center}
\indent Wollte man versuchen, allein an der Hand der von \textsc{Gauss} selbst durch
den Druck veröffentlichten Abhandlungen sich ein Bild zu entwerfen von den
Einsichten, die er sich auf dem Gebiete der Funktionentheorie\footnote{Wir bezeichnen mit diesem zu \textsc{Gauss}’ Zeiten weniger gebräuchlichen Terminus dasjenige wohlumgrenzte Gebiet der reinen Analysis, das man in neuerer Zeit mit diesem Namen belegt.} erworben
hat, so würde dieses Bild ein nach mehreren Seiten hin unvollständiges sein.
Hat doch \textsc{Gauss} vielleicht sparsamer als irgend ein anderer Forscher seinen
Zeitgenossen aus dem Schatze seiner Kenntnisse gespendet und überdies in
seinen Veröffentlichungen grösseren Wert auf eine vollendete Darstellung gelegt
als auf die Wiedergabe des Gedankenganges, der ihn zu seinen Resultaten
geführt hat. Dazu kommt noch die ihm eigentümliche Scheu, mit
Anschauungen und Methoden vor die Öffentlichkeit zu treten, die gar zu
umstürzlerisch erscheinen könnten, und das daraus entspringende Bestreben,
sich in seinen Publikationen den Gewohnheiten seiner Zeitgenossen anzupassen.
So schreibt er z. B. über seine Arbeiten zu den Grundlagen der Geometrie
am 27. Januar 1829 an \textsc{Bessel}\footnote{\textit{Briefwechsel Gauss-Bessel}, S. 490.} »Inzwischen werde ich wohl noch lange nicht
dazu kommen, meine \textls[120]{sehr ausgedehnten} Untersuchungen darüber zur öffentlichen
Bekanntmachung auszuarbeiten, und vielleicht wird dies auch bei meinen
Lebzeiten nie geschehen, da ich das Geschrei der Böotier scheue, wenn ich
meine Ansicht ganz aussprechen wollte.«
\HistoricalFoot{}{1\textsuperscript{*}}

\clearpage\SourceLeaf{89}\ArithmeticHead{4}{L}{L. SCHLESINGER, ÜBER GAUSS’ FUNKTIONENTHEORETISCHE ARBEITEN.}
Wir müssen darum neben den in Betracht kommenden von \textsc{Gauss} zum
Druck beförderten Arbeiten als mit diesen gleichberechtigte Quellen heranziehen:

1. Den Nachlass an handschriftlichen Aufzeichnungen wissenschaftlichen
Inhalts, der zur Zeit im Gaussarchiv der Göttinger Universitätsbibliothek aufbewahrt
wird und über dessen für uns in Betracht kommende Teile sogleich
berichtet werden soll\footnote{Vergl. die Bemerkungen über den Nachlass und die zur Zeit im Gaussarchiv angewandte Bezeichnungsweise der einzelnen Stücke, Werke X\textsubscript{1}, S. 576.}.

2. Den Briefwechsel, den \textsc{Gauss} mit zahlreichen Persönlichkeiten geführt
hat. — Bisher sind gesammelt im Druck erschienen die Briefwechsel mit
\textsc{Bessel}, \textsc{Wolfgang Bolyai}, \textsc{Lejeune Dirichlet}, \textsc{Gerling}, \textsc{Alexander von Humboldt},
\textsc{Nicolai}, \textsc{Olbers} und \textsc{Schumacher}\footnote{1. \textit{Briefwechsel zwischen Gauss und Bessel} [herausg. von A. \textsc{Auwers}], Leipzig 1880 (von 1804, Dez. 21 bis 1844, Aug. 15):
\par 2. \textit{Briefwechsel zwischen C. F. Gauss und W. Bolyai}, herausg. von F. \textsc{Schmidt} und P. \textsc{Stäckel}, Leipzig 1899 (von 1797, Sept. 29 bis 1853, Febr. 6);
\par 3. \textit{Briefwechsel zwischen C. F. Gauss und Chr. L. Gerling}, herausg. von \textsc{Clemens Schaefer}, Berlin 1927 (von 1810, Juni 26 bis 1859, Juni 30).
\par 4. \textit{Briefwechsel zwischen C. F. Gauss und P. G. Lejeune Dirichlet}, \textsc{Dirichlets} Werke II, Berlin 1897, S. 373—387 (von 1826, Mai 28 bis 1853, Februar 20).
\par 5. \textit{Briefe zwischen A. v. Humboldt und Gauss}, herausg. von K. \textsc{Bruhns}, Leipzig 1877 (von 1807, Juli 14 bis 1854, Dez. 4);
\par 6. \textit{Briefe von C. F. Gauss an B. Nicolai}, herausg. von W. \textsc{Valentiner}, Karlsruhe 1877 (von 1819, Aug. 11 bis 1845, Juni);
\par 7. \textit{Wilhelm Olbers, sein Leben und seine Werke}, herausg. von C. \textsc{Schilling}, II. Band, \textit{Briefwechsel zwischen Olbers und Gauss}, 1. Abt., Berlin 1900, 2. Abt., ebenda 1909 (von 1802, Jan. 18 bis 1839, Mai 30);
\par 8. \textit{Briefwechsel zwischen C. F. Gauss und H. C. Schumacher}, herausg. von C. A. F. \textsc{Peters}, Altona I (1860), II (1860), III (1861), IV (1862), V (1863), VI (1865) (von 1808, April 2 bis 1850, Nov. 4).}. Zahlreiche Briefe und Briefstellen
sind in den verschiedenen Bänden der Gaussschen Werke abgedruckt, insbesondere
enthält der Band XII (S. 228—324) Nachträge zu den gedruckten
Briefwechseln.

3. Nachrichten über mündliche Mitteilungen von \textsc{Gauss}, soweit sie uns
in mehr oder minder authentischer Form überliefert sind.

Der handschriftliche Nachlass ist schon von dem ersten Herausgeber der
Gaussschen Werke, E. J. \textsc{Schering}, in einem gewissen Masse verwertet worden.
Insbesondere hat \textsc{Schering} im Bande III der Werke eine Reihe von Abhand-\par
\clearpage\SourceLeaf{90}\ArithmeticHead{5}{R}{I. EINLEITENDES. DAS BENUTZTE MATERIAL.}
\noindent lungen zusammengestellt, die den Nachlassteilen entnommen sind, die sich
auf funktionentheoretische Gegenstände beziehen\footnote{1. \textit{Determinatio seriei nostrae}, S. 207—229; 2. \textit{De origine proprietatibusque generalibus numerorum arithmetico-geometricorum}, und \textit{De functionibus transcendentibus quae ex differentiatione mediorum arithmetico-geometricorum oriuntur}, S. 361—374; 3. \textit{Elegantiores integralis \(\displaystyle\int\frac{dx}{\sqrt{(1-x^4)}}\) proprietates}, S. 404—412; 4. \textit{De curva lemniscata}, S. 413—432; 5. \textit{Zur Theorie der neuen Transzendenten} I—V, S. 433—480; 6. \textit{Pentagramma mirificum}, S. 481—490. Wie \textsc{Schering} (Werke III, S. 492) berichtet, war die Herausgabe des auf die Theorie der elliptischen Funktionen bezüglichen Nachlasses von der Gesellschaft der Wissenschaften \textsc{Riemann} übertragen worden, der (vergl. \textsc{Hattendorf}, \textit{Die elliptischen Funktionen in dem Nachlasse von Gauss}, Hannover 1869) sich auch schon mit dem Studium der Handschriften beschäftigt hatte, als er durch Krankheit und die dadurch bedingte Abwesenheit von Göttingen genötigt war, diese Arbeiten zurückzustellen. — Nach seinem Tode (1866, Juli) übernahm \textsc{Schering} die Handschriften; der Band III der Werke, der damals bis auf den die elliptischen Funktionen behandelnden Teil bereits gedruckt war (daraus erklärt sich die Jahreszahl 1866 auf dem Titelblatt), wurde dann im November 1868 abgeschlossen (vergl. die Angaben in E. \textsc{Scherings} Mathem. Werken, II, 1906, S. 463).}. Ein Teil dieser Abhandlungen
ist eine im wesentlichen unveränderte Wiedergabe der Gaussschen
Handschriften; \textsc{Schering} hat aber oft die von ihm abgedruckten Aufzeichnungen,
ohne Rücksicht auf die Zeit ihrer Entstehung, nach Gegenständen
geordnet, wohl weil es ihm unmöglich war, viele dieser Aufzeichnungen auch
nur annähernd zu datieren. Über die Herkunft der einzelnen Stücke geben
die Werke III, S. 491—496 abgedruckten Bemerkungen \textsc{Scherings} einige
Hinweise. Die Abhandlung \textit{Fortsetzung der Untersuchungen über das arithmetisch-geometrische Mittel}
(Werke III, S. 375—402) kann als ein von \textsc{Scherings}
Hand aufgeführtes Bauwerk bezeichnet werden, in das einzelne von \textsc{Gauss} herrührende
Steine eingefügt sind. Sie kann demnach als Material für die Wiederherstellung
des Gaussschen Gedankenganges nicht eigentlich benutzt werden.

Eine Nachlese des funktionentheoretischen Nachlasses hat \textsc{Fricke} (1900) im
VIII. Bande der Werke gegeben\footnote{Werke VIII, S. 65—67; 69—75; 76—79; 80—83; 84—85; 93—94; 96; 98; 99—105; 106—111.}. \textsc{Fricke} war schon in der Lage, das von \textsc{Stäckel}
1898 bei \textsc{Gauss}’ Nachkommen aufgefundene \textit{Tagebuch} (oder Notizenjournal)
zu benutzen, in das \textsc{Gauss} von 1796 bis 1814 kurze Notizen in lateinischer
Sprache über die Gegenstände seiner wissenschaftlichen Beschäftigung, zumeist
mit genauer Datierung einzutragen pflegte. Dieses \textit{Tagebuch} wurde 1901 von
\textsc{Felix Klein} in der \textit{Festschrift zur Feier des hundertfünfzigjährigen Bestehens
der Königlichen Gesellschaft der Wissenschaften zu Göttingen}\footnote{\textit{Beiträge zur Gelehrtengeschichte Göttingens}, Berlin 1901, S. 1; abgedruckt Mathem. Annalen 57, 1903, S. 1. — Eine neue Ausgabe mit ausführlichen Erläuterungen und einem Faksimile der Handschrift befindet sich Werke X\textsubscript{1}, S. 483—574.} herausgegeben und\par
\clearpage\SourceLeaf{91}\ArithmeticHead{6}{L}{L. SCHLESINGER, ÜBER GAUSS’ FUNKTIONENTHEORETISCHE ARBEITEN.}
\noindent mit Anmerkungen versehen, in denen auch eine Reihe anderer Aufzeichnungen
von \textsc{Gauss} zum ersten Male veröffentlicht worden ist. In diesen Anmerkungen
hat \textsc{Klein} auch zuerst auf die Wichtigkeit der Notizen hingewiesen, die \textsc{Gauss}
auf die Durchschussblätter eines ihm gehörigen Buches eingetragen hat, das
wir kurz mit dem Namen seines Verfassers als den \textit{Leiste} bezeichnen. Es ist
nämlich ein Exemplar von \textsc{Chr. Leistes} 1790 in Wolfenbüttel erschienenem
Werkchen: \textit{Die Arithmetik und Algebra}, das \textsc{Gauss} sich angeschafft, das er hat
binden und mit Schreibpapier durchschiessen lassen\footnote{Neben seinem Namen hat \textsc{Gauss} auf das Schutzblatt des Buches eingeschrieben:
\begin{center}\begin{tabular}{r@{\quad}r}
Const. liber ipse&8\\
ligatura&4\\\cline{2-2}
&12
\end{tabular}\end{center}
\noindent Dadurch wird bestätigt, dass wirklich alle handschriftlichen Eintragungen, die das Buch enthält, von \textsc{Gauss} herrühren. — In Bezug auf den \textit{Leiste} schrieb F. \textsc{Klein} nach dem Empfange des Manuskripts der ersten Bearbeitung dieses Aufsatzes an den Verfasser aus Hahnenklee am 12. Januar 1912 das folgende: »Ich habe jetzt .\ .\ . Ihren Essai etwas genauer durchgesehen und will Ihnen vor allen Dingen aussprechen, dass mir der Essai nach Anlage und Ausführung ganz besonderes Vergnügen macht. Ich selbst habe diese Dinge nie im Zusammenhange durchgearbeitet, sondern nur so viel, als mir 1901, wo ich das \textit{Tagebuch} als Gelegenheitsschrift herausgeben wollte, gerade in die Hände fiel. Die Einsicht, dass im \textit{Leiste} ganz besonders wichtige Dinge stehen, kam mir als volle Überraschung erst während der Korrektur.«}. Auf den Durchschussblättern
stehen zunächst einige zerstreute Notizen, die sich auf den Text des
Buches beziehen, eine davon ist Werke X\textsubscript{1}, S. 78 abgedruckt; diese Notizen
stammen sicher noch aus der Braunschweiger Zeit vor dem Abgange an die
Universität Göttingen (11. Oktober 1795). Dann hat \textsc{Gauss} auf den ersten
Durchschussblättern ein Verzeichnis der Schriften \textsc{Eulers} (nach der Fussschen
Liste von 1783) eingetragen, wohl in seinem ersten Göttinger Studiensemester, und
weiterhin hat ihm der \textit{Leiste} geradezu als das erste einer langen Reihe von
Notizbüchern gedient, auf deren Blättern er während seiner fast sechzigjährigen
wissenschaftlichen Tätigkeit den in ihm emporkeimenden Gedanken die erste
literarische Form gegeben hat. Auf den Durchschussblättern des \textit{Leiste} spiegeln
sich auf diese Weise gewissermassen die einzelnen Entwicklungsstufen des
jugendlichen \textsc{Gauss}; wir sehen den Schüler des Braunschweiger Carolinums,
der den \textit{Leiste} studiert und sich Notizen zum Texte macht, dann den Göttinger
Studenten, der \textsc{Euler} zu seinem Lehrmeister macht, »dessen Studium doch
stets die beste durch nichts anderes zu ersetzende Schule für die verschiedenen
mathematischen Gebiete bleiben wird«\footnote{Aus einem Briefe von \textsc{Gauss} an P. H. v. \textsc{Fuss} vom 16. September 1849.}, endlich den zum frei schaffenden\par
\clearpage\SourceLeaf{92}\ArithmeticHead{7}{R}{I. EINLEITENDES. DAS BENUTZTE MATERIAL.}
\noindent Mathematiker gewordenen \textsc{Gauss}, der alsbald nicht nur in der Zahlentheorie
und Algebra, sondern auch in der Analysis auf selbstgebahnten Wegen vorwärts
dringt. Die grosse Bedeutung, die neben dem \textit{Tagebuch} dem \textit{Leiste} für
die Erforschung von \textsc{Gauss}’ wissenschaftlicher Entwicklung in den ersten Jahren
seiner Laufbahn zukommt, liegt danach auf der Hand. \textsc{Schering} hat weder
das \textit{Tagebuch} noch den \textit{Leiste} gekannt, daraus erklärt sich die oben bemerkte
Unmöglichkeit für ihn, viele der von ihm bearbeiteten Handschriften zu datieren.
Man kann wohl sagen, dass durch die Auffindung dieser beiden
wichtigen Dokumente die Gaussforschung neu belebt worden ist, indem es
erst an ihrer Hand möglich wurde, Einsicht zu gewinnen in den Entwicklungsgang
gerade der funktionentheoretischen Arbeiten von \textsc{Gauss}. — Ermöglichen
doch die Eintragungen des \textit{Tagebuchs} nicht nur die Datierung der in verschiedenen
Heften und Handbüchern sowie auf Zetteln befindlichen Aufzeichnungen,
sondern geben auch Kunde von Resultaten, von denen in dem übrigen
handschriftlichen Nachlass keine Spur vorhanden ist.

Die Vergleichung hat ergeben, dass eine Anzahl von in fortlaufender
Folge eingetragenen Leistenotizen fast lückenlos mit Aufzeichnungen des \textit{Tagebuchs}
aus den Jahren 1796—1798 korrespondiert; \textsc{Gauss} hat also in diesen
für seine Entwicklung so wichtigen Jahren, den \textit{Leiste} \textls[120]{regelmässig} zu
wissenschaftlichen Aufzeichnungen benutzt. Neben diesen Aufzeichnungen,
die wir als \textls[120]{reguläre} bezeichnen, finden sich aber solche, die an früher frei
gebliebenen Stellen, an den andersartigen Schriftzügen und anderer Tinte meist
deutlich erkennbar, später eingetragen worden sind. Für die Datierung dieser
\textls[120]{irregulären} Leistenotizen bietet also die Reihenfolge der Blätter, auf denen
sie aufgezeichnet sind, keinen Anhaltspunkt. — Die von \textsc{Klein} in den Anmerkungen
zur ersten Ausgabe des \textit{Tagebuchs} wiedergegebenen Leistenotizen
gehören zu den regulären, ein Teil der irregulären wurde 1912 im II. Hefte
der \textit{Materialien}\footnote{\textit{Materialien für eine wissenschaftliche Biographie von Gauss}, gesammelt von F. \textsc{Klein} und M. \textsc{Brendel}, II, C. F. \textsc{Gauss}, \textit{Fragmente zur Theorie des arithmetisch-geometrischen Mittels aus den Jahren 1797—1799}, herausgegeben und erläutert von L. \textsc{Schlesinger}.} veröffentlicht. Dieses Heft brachte auch noch einige andere,
auf die Theorie der Modulfunktion bezügliche Aufzeichnungen aus dem als
Scheda Ac bezeichneten Notizheftchen (begonnen November 1799), die mit
Tagebuchnotizen der Jahre 1799—1800 korrespondieren und aus denen einige\par
\clearpage\SourceLeaf{93}\ArithmeticHead{8}{L}{L. SCHLESINGER, ÜBER GAUSS FUNKTIONENTHEORETISCHE ARBEITEN.}
\noindent Formeln schon in \textsc{Kleins} Anmerkungen zum \textit{Tagebuch} angeführt worden
waren. Es folgte im III. Heft der \textit{Materialien} (ebenfalls 1912), das die erste
Bearbeitung des vorliegenden Aufsatzes enthielt, noch eine Reihe weiterer
Nachlassstücke. Alle diese Stücke liegen nunmehr mit noch zahlreichen andern
bis dahin unveröffentlichten Aufzeichnungen des Nachlasses planmässig
geordnet und erläutert in dem 1917 erschienenen Bande X, 1 der Werke vor\footnote{Werke X, 1, Abschnitt Analysis, S. 138—449.},
einige kleinere Stücke sind den Erläuterungen zum \textit{Tagebuch} eingefügt worden
(Werke X, 1, S. 485ff.). Eine kleine Nachlese bringt der Abschnitt Varia,
Werke XII (1929)\footnote{Wir geben hier eine Aufzählung der Nachlassteile, die für \textsc{Gauss}’ funktionentheoretische Arbeiten hauptsächlich in Betracht kommen; die den einzelnen Stücken vorgesetzten Zeichen sind die Signaturen, unter denen die Stücke im Gaussarchiv zu finden sind.
\begin{list}{}{\setlength{\leftmargin}{1.7em}\setlength{\rightmargin}{0pt}\setlength{\itemsep}{0pt}\setlength{\parsep}{0pt}}
\item[1.] Das \textit{Tagebuch} (Notizenjournal).
\item[2.] Die Schedae.
\begin{list}{}{\setlength{\leftmargin}{1.4em}\setlength{\itemsep}{0pt}\setlength{\parsep}{0pt}}
\item Aa: Schedae Nr. I. Jul. 1798.
\item Ab: Exercitationes atque Schedae analyticae, 1798 Nov.
\item Ac: Varia, Novbr. 1799. Imprimis de Integrali \(\displaystyle\int\frac{d\varphi}{\sqrt{(1+\mu\mu\sin\varphi^2)}}\mathpunct{\vcenter{\hbox{.}}}\)
\item Ae: Varia, Jul. 1800.
\item Af: Mémoires de Mathématiques, Bronsuic 1801.
\item An: Cereri, Palladi, Junoni sacrum, Febr. 1805.
\end{list}
\item[3.] \textsc{Leiste} »Die Arithmetik und Algebra«, durchschossenes Exemplar mit handschriftlichen Aufzeichnungen von \textsc{Gauss}. (Im folgenden kurz mit \textit{Leiste} zitiert).
\item[4.] Die Handbücher:
\begin{list}{}{\setlength{\leftmargin}{1.4em}\setlength{\itemsep}{0pt}\setlength{\parsep}{0pt}}
\item 15 (Ba). Opuscula varii Argumenti. Volumen Primum. Brunovici 1800.
\item 16 (Bb). Den astronomischen Wissenschaften gewidmet. November 1801.
\item 17 (Bc). Astronomische Untersuchungen und Rechnungen vornehmlich über die Ceres Ferdinandea. 1802.
\item 18 (Bd). Mathematische Brouillons. Oktober 1805.
\item 19 (Be). Kleine Aufsätze aus verschiedenen Teilen der Mathematik. Angefangen im Mai 1809.
\end{list}
\item[5.] Zwei Kapseln mit Zetteln, und auf einzelne Blätter geschriebenen grösseren Entwürfen, die in einzelne Umschläge geordnet sind. (Die folgenden Überschriften sind nicht von \textsc{Gauss}.)
\begin{list}{}{\setlength{\leftmargin}{1.4em}\setlength{\itemsep}{0pt}\setlength{\parsep}{0pt}}
\item Fa. Reihen im allgemeinen.
\item Fb. \textsc{Lagrangesche} Reihe.
\item Fc. Reihe \(1+\dfrac{\alpha\cdot\beta}{1\cdot\gamma}x+{}\) etc.
\item Fd. Tafeln der \(\Pi\)-Funktion.
\item Ff. Arithm.-geom. Mittel, ältere Untersuchungen.
\item Fg. Modulfunktion.
\item Fh. Lemniskatische Funktionen.
\item Fi. Allgemeine elliptische Funktionen und Integrale.
\item Fm. Varia analytica.
\end{list}
\end{list}}.

\clearpage\SourceLeaf{94}\ArithmeticHead{9}{R}{II. DIE ERSTE JUGENDZEIT BIS OKTOBER 1795.}
Obwohl nun dieses reichhaltige Material im Ganzen einen befriedigenden
Einblick in den Entwicklungsgang von \textsc{Gauss}’ funktionentheoretischem Schaffen
gewährt, bleiben doch im einzelnen noch manche Lücken, die durch Vermutungen
ausgefüllt werden müssen, wie ja auch sonst, der Natur der Sache
nach, der subjektiven Auffassung ein gewisser Raum verstattet werden muss.

In mannigfacher Weise wird der Bericht, den wir über \textsc{Gauss}’ Tätigkeit
gerade in den entscheidenden Jahren zu erstatten haben werden, durch die
anderen Abhandlungen dieses Bandes X, 2, sowie durch die des Bandes XI, 2
ergänzt; besonders stehen die funktionentheoretischen Untersuchungen von
\textsc{Gauss} in inniger Wechselbeziehung zu seinen Arbeiten über Zahlentheorie,
Geometrie und Astronomie.

Es möge noch hervorgehoben werden, dass wir es hier ausschliesslich mit
\textsc{Gauss} zu tun haben; nicht um eine Entwicklungsgeschichte der Analysis im
XIX. Jahrhundert handelt es sich, sondern einzig und allein um eine Geschichte
der Entwicklung von \textsc{Gauss}’ Gedanken. Dem Einfluss, den die von
\textsc{Gauss} nicht veröffentlichten Untersuchungen mittelbar oder unmittelbar auf
die Weiterentwicklung der Funktionentheorie ausgeübt haben, werden wir von
Fall zu Fall nachzugehen bestrebt sein; hier möge nur noch betont werden,
dass es nicht beabsichtigt ist, aus dem von \textsc{Gauss} bei seinen Lebzeiten zurückgehaltenen
Material Prioritätsansprüche zu Gunsten von \textsc{Gauss} gegenüber
solchen Mathematikern zu konstruieren, die Resultate, die \textsc{Gauss} für sich
schon früher entwickelt, aber nicht bekannt gemacht hatte, später oder gleichzeitig
durchaus unabhängig und von diesen Untersuchungen von \textsc{Gauss} unbeeinflusst
gefunden haben. Die Feststellung, dass ein Mathematiker eine Entdeckung
gemacht hat, die auch dem Geiste des princeps mathematicorum
entsprungen war, soll den Ruhm des zweiten Entdeckers nicht schmälern,
sondern ihn erhöhen!
